\subsection{Assigning Coordinates by Orthogonal Projection}

Only after the axes, directions, common unit, integer marks, refined scale, and
real labels have been constructed can an arbitrary point of the plane be given
coordinates.

Let $P$ be any point of the plane. Through $P$, construct the line perpendicular
to the $x$-axis. It meets the $x$-axis at a unique point $P_x$. Through $P$,
construct the line perpendicular to the $y$-axis. It meets the $y$-axis at a
unique point $P_y$.

\begin{center}
\begin{tikzpicture}[scale=0.9]
    \draw[axis] (-1.1,0) -- (5.0,0) node[right] {$x$};
    \draw[axis] (0,-1.1) -- (0,4.0) node[above] {$y$};
    \coordinate (P) at (3.2,2.4);
    \coordinate (Px) at (3.2,0);
    \coordinate (Py) at (0,2.4);
    \draw[helper] (P) -- (Px);
    \draw[helper] (P) -- (Py);
    \fill (P) circle (2.2pt) node[above right] {$P$};
    \fill (Px) circle (1.8pt) node[below] {$P_x$};
    \fill (Py) circle (1.8pt) node[left] {$P_y$};
\end{tikzpicture}
\end{center}

If the signed real labels of $P_x$ and $P_y$ are $x$ and $y$, respectively, we
assign to $P$ the ordered pair
\[
P\longleftrightarrow(x,y).
\]

The construction is unambiguous because the perpendicular through a given point
to a given line is unique. Conversely, given real numbers $x$ and $y$, locate
the corresponding points on the two axes and construct through them lines
perpendicular to their respective axes. Those two lines meet in exactly one
point of the plane.

\begin{theorembox}
Once the two perpendicular axes, their positive directions, and a common unit
are fixed, orthogonal projection gives a one-to-one correspondence
\[
\boxed{
\text{points of the Euclidean plane}
\longleftrightarrow
\text{pairs }(x,y)\in\mathbb R^2
}.
\]
The ordered pair is the numerical address of a geometric point in the chosen
coordinate system. It is not the point itself and it is not the source of the
geometry.
\end{theorembox}
